\label{chap.7}

This chapter reports three evaluations of the hybrid tester. The pre-registered held-out
evaluation on the validation suite (Section~\ref{sec:heldout}) is complete. The final tuning
round on the rebuilt development suite (Section~\ref{sec:tuning}) and the external comparison
with WebExplor and WebRLED (Section~\ref{sec:extresults}) follow it.

\section{Held-out evaluation (pre-registered)}
\label{sec:heldout}

\subsection{Design}

The pre-registration was frozen on the project owner's approval, and the commit was tagged,
before any agent ran on any validation application.

\begin{itemize}
  \item \textbf{Applications}: the eight validation applications (Table~\ref{tab:valsuite}),
        24 seeded faults.
  \item \textbf{Arms}:
        \begin{itemize}[noitemsep]
          \item \code{random}, uniform over the page's valid actions;
          \item \code{hybrid\_nolearn}, the primary arm;
          \item \code{hybrid\_bandit};
          \item \code{v1\_dqn\_pretrained}, v1's best AC-DQN arm from Experiment 1: pretrained
                without the LLM reward and adapting online, using that experiment's own
                checkpoints. Seed $s$ loads the checkpoint of fold $s \bmod 6$, pretraining
                seed $s \bmod 3$.
        \end{itemize}
  \item \textbf{Runs}: seeds 0--4, $B = 2{,}000$ interactions, 160 hunts, plus 32 clean-build
        control runs (one per arm and application). The run took about seven and a half
        hours with four workers and no other workload.
  \item \textbf{Hypotheses} (difference in restricted mean $T$, decision rule of
        Section~\ref{sec:stats}):
        \begin{itemize}[noitemsep]
          \item H1 (primary): the hybrid tester is faster than random;
          \item H2: the bandit makes the hybrid faster;
          \item H3: the hybrid with the bandit is faster than random;
          \item H4: the hybrid is faster than v1's best AC-DQN.
        \end{itemize}
\end{itemize}

One of 192 jobs failed for an infrastructure reason: Windows held a file lock while the
temporary build folder of a clean-build control was being deleted. As pre-registered, the job
was rerun unchanged; it completed with no finding.

\subsection{Results}

\begin{table}[H]
\centering
\small
\caption{Held-out evaluation: pre-registered comparisons (negative $\Delta$ = first arm faster).}
\label{tab:heldout-h}
\resizebox{\textwidth}{!}{%
\begin{tabular}{@{}llrrrrl@{}}
\toprule
& Comparison & Median $\Delta$ & 95\% interval & Apps favouring & Log-rank $p$ & Verdict \\
\midrule
H1 & \code{hybrid\_nolearn} $-$ \code{random} & $-324$ & $[-665, +148]$ & 6/8 & 0.159 & inconclusive \\
H2 & \code{hybrid\_bandit} $-$ \code{hybrid\_nolearn} & $+86$ & $[-304, +285]$ & 4/8 & 0.781 & inconclusive \\
H3 & \code{hybrid\_bandit} $-$ \code{random} & $-191$ & $[-801, +67]$ & 5/8 & 0.180 & inconclusive \\
H4 & \code{hybrid\_nolearn} $-$ \code{v1\_dqn\_pretrained} & $+42$ & $[-196, +574]$ & 4/8 & 0.276 & inconclusive \\
\bottomrule
\end{tabular}}
\end{table}

\begin{table}[H]
\centering
\small
\caption{Held-out evaluation per arm (40 runs each).}
\label{tab:heldout-arms}
\resizebox{\textwidth}{!}{%
\begin{tabular}{@{}lrrrrr@{}}
\toprule
Arm & Found a fault & RMST & Faults per run & Distinct faults & Clean false positives \\
\midrule
\code{v1\_dqn\_pretrained} & 34 (0.85) & \textbf{595} & 1.18 & 17 / 24 & 0 / 8 \\
\code{hybrid\_nolearn} & 30 (0.75) & 737 & 1.15 & 18 / 24 & 0 / 8 \\
\code{hybrid\_bandit} & 34 (0.85) & 769 & \textbf{1.23} & \textbf{19} / 24 & 0 / 8 \\
\code{random} & 32 (0.80) & 978 & 1.00 & 13 / 24 & 0 / 8 \\
\bottomrule
\end{tabular}}
\end{table}

\begin{table}[H]
\centering
\small
\caption{Held-out evaluation per application: runs that found a fault (of 5) and mean $T$.}
\label{tab:heldout-apps}
\begin{tabular}{@{}lllll@{}}
\toprule
App & \code{random} & \code{v1\_dqn\_pretrained} & \code{hybrid\_nolearn} & \code{hybrid\_bandit} \\
\midrule
dutyday & 4/5 $\cdot$ 1{,}217 & \textbf{5/5 $\cdot$ 281} & 4/5 $\cdot$ 788 & 4/5 $\cdot$ 778 \\
fosterly & 5/5 $\cdot$ 1{,}262 & 4/5 $\cdot$ 763 & 4/5 $\cdot$ 681 & \textbf{5/5 $\cdot$ 370} \\
larder & 5/5 $\cdot$ 411 & 5/5 $\cdot$ 407 & \textbf{5/5 $\cdot$ 211} & 5/5 $\cdot$ 393 \\
markbook & 1/5 $\cdot$ 1{,}805 & \textbf{4/5 $\cdot$ 567} & 3/5 $\cdot$ 1{,}141 & 4/5 $\cdot$ 1{,}442 \\
quizcraft & 5/5 $\cdot$ 299 & 5/5 $\cdot$ 97 & \textbf{5/5 $\cdot$ 81} & 5/5 $\cdot$ 366 \\
sideline & 2/5 $\cdot$ 1{,}623 & 1/5 $\cdot$ 1{,}759 & \textbf{4/5 $\cdot$ 620} & 3/5 $\cdot$ 822 \\
skylark & 5/5 $\cdot$ 980 & \textbf{5/5 $\cdot$ 614} & 0/5 $\cdot$ 2{,}000 & 3/5 $\cdot$ 1{,}696 \\
tally & 5/5 $\cdot$ \textbf{229} & 5/5 $\cdot$ 276 & 5/5 $\cdot$ 376 & 5/5 $\cdot$ 284 \\
\bottomrule
\end{tabular}
\end{table}

All four pre-registered comparisons are inconclusive (Table~\ref{tab:heldout-h}). As the
pre-registration anticipated, eight applications give the median rule wide intervals, and no
comparison was supported or refuted. The per-fault counts are in
Appendix~\ref{app:faults}.

\subsection{What the held-out data shows (descriptive)}

\begin{enumerate}
  \item \textbf{The held-out faults were easier for random exploration than the development
        ones.} Random verified a fault in 80\% of runs, against 43\% (18 of 42) on the
        development suite. Several faults sit one or two interactions from a page that
        random visits constantly. Three faults were verified in 5 of 5 runs by every arm,
        which compresses every difference.
  \item \textbf{The hybrid arms reached the deeper faults.} 22 of the 24 faults were
        verified by at least one arm; two by none. Of the 16 faults random verified in at
        most one of five runs, the two hybrid arms together verified 14 at least once, and
        \code{v1\_dqn\_pretrained} 9. Nine faults were verified by a hybrid arm and by no
        random run. This is where test operators help: completing forms validly, invalid
        input, edits after review.
  \item \textbf{One clear blind spot: list paging.} All three skylark faults sit behind its
        feed's "Load more" button. Random presses it often by chance. The hybrid's operators
        work on forms and none pages through a list, so \code{hybrid\_nolearn} reached the
        right pages in every run but never repeated the paging (0 of 5). This cost H1 one
        of its eight applications. It was reported, not fixed, because the code was frozen;
        it motivated the \code{browse\_list} operator (Section~\ref{sec:v3-options}).
  \item \textbf{v1's pretrained AC-DQN did better here than on its own benchmark}: discovery
        0.85 and the lowest mean time, notably on dutyday and markbook. Whether that
        reflects the method or the applications' shallower faults cannot be separated with
        eight applications (H4 inconclusive).
  \item \textbf{The learner (H2).} \code{hybrid\_bandit} found a fault in more runs than
        \code{hybrid\_nolearn} (34 against 30) and more distinct faults (19 against 18), with
        a slightly higher mean time (769 against 737). As in development, the learner is
        not distinguishable from no learner.
  \item \textbf{Zero false positives.} No clean-build control verified a seeded fault
        (0 of 32 runs).
\end{enumerate}

\paragraph{Specification-free invariants.} As pre-registered, they are reported separately.
They flagged no seeded fault for any arm. On clean builds they produced three reproduced
candidates (two for \code{hybrid\_nolearn}, one for \code{hybrid\_bandit}), all on
quizcraft's editor, all of the form "an entered value is gone after reload". On inspection,
quizcraft keeps its working copy in page memory until the user saves, so losing unsaved
edits on reload is intended. They are false alarms: the reload invariant cannot tell saved
state from unsaved state.

\section{Tuning round v3 on the rebuilt development suite}
\label{sec:tuning}

\subsection{Design}

After the held-out evaluation, the development suite was rebuilt (Section~\ref{sec:devsuite})
and the tester re-tuned on it before the external comparison. Eight arms ran on the eight
development applications with seeds 0--2, 192 runs of 2{,}000 interactions:

\begin{itemize}[noitemsep]
  \item the references: \code{random}, \code{hybrid\_nolearn}, \code{hybrid\_bandit} and
        \code{v1\_dqn\_pretrained};
  \item one option each (Section~\ref{sec:v3-options}):
        \begin{itemize}[noitemsep]
          \item \code{dev\_v3\_browse},
          \item \code{dev\_v3\_clickables},
          \item \code{dev\_v3\_curiosity},
          \item \code{dev\_v3\_failfeedback}.
        \end{itemize}
\end{itemize}

Each run records verified faults (including crash-type faults), unique failures and
JavaScript coverage. These are development data: they choose the configuration to freeze,
and are not evidence of generalisation.

\subsection{Results}

All 192 runs completed, none failed. Table~\ref{tab:tuning-arms} summarises each arm.
Table~\ref{tab:tuning-matched} compares each arm with \code{hybrid\_nolearn} on matched
(application, seed) pairs.

\begin{table}[H]
\centering
\small
\caption{Tuning round v3 per arm: 24 runs each on the eight rebuilt development
applications, which hold 41 seeded faults (8 of them crash-type).}
\label{tab:tuning-arms}
\resizebox{\textwidth}{!}{%
\begin{tabular}{@{}lrrrrrr@{}}
\toprule
Arm & Found a fault & RMST & Distinct faults & Crash faults & Unique failures / run & JS coverage \\
\midrule
\code{dev\_v3\_failfeedback} & 21 / 24 & \textbf{450} & 19 / 41 & 3 / 8 & 3.79 & 0.893 \\
\code{random} & \textbf{23} / 24 & 498 & 19 / 41 & 1 / 8 & 5.38 & 0.861 \\
\code{dev\_v3\_browse} & 20 / 24 & 545 & 19 / 41 & 3 / 8 & 3.12 & 0.885 \\
\code{hybrid\_nolearn} & 21 / 24 & 560 & 19 / 41 & \textbf{4} / 8 & 2.88 & 0.896 \\
\code{v1\_dqn\_pretrained} & 20 / 24 & 564 & 19 / 41 & 2 / 8 & 4.75 & \textbf{0.909} \\
\code{dev\_v3\_clickables} & 19 / 24 & 638 & 17 / 41 & 3 / 8 & 2.50 & 0.891 \\
\code{dev\_v3\_curiosity} & 20 / 24 & 650 & 16 / 41 & 2 / 8 & \textbf{5.67} & 0.894 \\
\code{hybrid\_bandit} & 19 / 24 & 692 & 16 / 41 & 3 / 8 & 2.83 & 0.896 \\
\bottomrule
\end{tabular}}
\end{table}

\begin{table}[H]
\centering
\small
\caption{Tuning round v3 against \code{hybrid\_nolearn}: per-application mean difference in
$T$ (negative = faster), median over applications with a 95\% bootstrap interval. Apps
with no difference are not counted as faster or slower.}
\label{tab:tuning-matched}
\begin{tabular}{@{}lrrc@{}}
\toprule
Arm & Median & 95\% interval & Apps faster / slower \\
\midrule
\code{dev\_v3\_browse} & $-20$ & $[-79, 0]$ & 6 / 1 \\
\code{dev\_v3\_failfeedback} & 0 & $[-379, 0]$ & 3 / 1 \\
\code{random} & $-18$ & $[-285, +186]$ & 5 / 3 \\
\code{v1\_dqn\_pretrained} & $-45$ & $[-182, +265]$ & 5 / 3 \\
\code{dev\_v3\_clickables} & 0 & $[0, +286]$ & 0 / 3 \\
\code{dev\_v3\_curiosity} & $+43$ & $[-248, +549]$ & 4 / 4 \\
\code{hybrid\_bandit} & $+140$ & $[+112, +416]$ & 1 / 7 \\
\bottomrule
\end{tabular}
\end{table}

\paragraph{Reading.}
\begin{itemize}
  \item \textbf{\code{browse\_list} helps consistently but modestly.} It was faster on six of
        seven applications and slower on one, and lost no fault.
  \item \textbf{Failure feedback gave the best restricted mean time} (450).
  \item \textbf{Custom clickables did not help.} They were never faster, and verified fewer
        distinct faults.
  \item \textbf{The bandit was slower than no learner} on seven of eight applications. Its
        interval excludes zero.
  \item \textbf{Curiosity Q-learning triggered the most unique browser errors} (5.67 per
        run) but verified fewer seeded faults.
  \item \textbf{Random exploration was strong on these applications.} It found a fault in 23
        of 24 runs and triggered many errors, but verified only one crash-type fault.
  \item \textbf{The rebuilt applications are hard.} 15 of the 41 faults were verified by no
        arm.
\end{itemize}

\subsection{Configuration chosen}

The two options that did not hurt any measure, \code{browse\_list} and failure feedback,
were combined (\code{dev\_v3\_combined}) and checked in 24 further runs on the same
applications and seeds before freezing. The combination found a fault in 22 of 24 runs,
with restricted mean $T$ 465. It verified 22 of the 41 faults, more than any single arm, and
4 of the 8 crash-type faults. Against \code{hybrid\_nolearn} it was faster on six of seven
applications (median $-52$, interval $[-145, 0]$). It was frozen as \code{hybrid\_final},
and the commit was tagged before anyone looked at the external benchmark's applications.
Custom clickables and the bandit were not adopted.

\section{Comparison with WebExplor and WebRLED}
\label{sec:extresults}

\subsection{Design}

The WebRLED paper~\cite{gu2025webrled} reports results on WebExplor's six subject
applications (Dimeshift, Pagekit, Splittypie, Phoenix, Retroboard and Petclinic), with
one-hour runs and five repetitions. WebExplor averaged 55.54\% client-side branch coverage
and 3.1 failures; WebRLED averaged 64.18\% and 4.1. Coverage was measured with Istanbul on
the instrumented client code. Failures were console errors, deduplicated by hand.

The comparison was pre-registered when \code{hybrid\_final} was frozen, before anyone
looked at the subject applications.
\begin{itemize}
  \item \textbf{Applications.} The same six applications, from the Docker images of the
        DIG replication material~\cite{biagiola2019diversity}, which WebExplor's benchmark
        reuses. Each image instruments its client JavaScript with Istanbul before the
        application starts.
  \item \textbf{Arms.} \code{hybrid\_final}; the hybrid with curiosity Q-learning, the RL
        arm closest to WebExplor and WebRLED; and random.
  \item \textbf{Runs.} One hour of wall-clock time per run, three repetitions (seeds 0--2),
        and a fresh application instance per run, so $3 \times 6 \times 3 = 54$ runs.
  \item \textbf{Measures.} Branch coverage of the instrumented client code, read from
        Istanbul's counters after every step and unioned over the run; and unique failures,
        meaning console errors, page errors and failed requests, deduplicated automatically
        by kind, page and normalised message.
  \item \textbf{Analysis.} A descriptive comparison with the published averages, with no
        significance test against numbers from other machines and runs, and a per-application
        comparison between our arms.
\end{itemize}

Two decisions were made during set-up and recorded before any run. First, Pagekit and
Phoenix start with a sign-in macro using their documented accounts. Second, requests to a
development live-reload server and to web analytics, which the images' pages make, are
answered with an empty response. Left alone, they would add the same console failure to
every run of every arm. Two differences from the published protocol remain and are
reported: three repetitions instead of five, and automatic instead of manual deduplication
of failures.

\subsection{Results}

\pending{Coverage and failures per application against the published numbers; verified
findings.}

\section{Discussion}
\label{sec:discussion}

\pending{Discussion across all evaluations: what the structure contributes, what learning
contributes, how the tester compares with the literature, and the limits of the evidence.}
